% 本文はsectionを最大単位にしているため、付録もsectionで採番する。
\appendix
\setcounter{section}{0}
\renewcommand{\thesection}{\Alph{section}}

\section{本文を補う資料}
\label{sec:appendix}
\begin{itemize}
    \item 数式の詳しい導出，実験の設定一覧，障害の種類ごとの詳しい結果を置く．本文は研究の問いと答えの流れを追える長さに保つ．
    \item 採用しなかった案は，本文の議論を支えるものだけを補足する．開発履歴をすべて並べる必要はない．
\end{itemize}

% ----------------------------------------------------------------------
% 以下は執筆者用の根拠対応表。レビュー用PDFには表示しない。
% 2026-09-01確認。数値を本文へ記入する際は、元結果と設定の対応を再確認する。
% draft1で段落・具体例・引用へ展開し、draft2で設定と根拠を確定する。
% draft3で主張と限界を精査し、draft4--5で表現・体裁・提出要件を確認する。
%
% E1：提案手法（序論・背景・提案手法）
% src/models/amber.py
% configs/main/rcaeval_re1_zenodo_v2.yaml
% notes/decisions/2026-08-28_final_amber.md
% 次数はAR(3)が現行設定。最終決定についてはE6の後続判断を優先する。
%
% E2：データと評価条件（背景・評価実験）
% notes/decisions/2026-08-19_rcaeval_re1_zenodo_v2_migration.md
% notes/decisions/2026-08-25_rcaeval_input_column_sanitation.md
% src/evaluation/metrics.py
% RCAEval RE1 Zenodo record 14590730 v2、各125件、計375件、前後最大600点。
% 前後結合区間による完全定数列の除外と、正常限定のモデル学習を区別する。
% 旧データからの移行経緯はdraft0の読者向け説明には含めない。
%
% E3：現在の主方式の結果・ケース対応診断（問い1・問い2・考察）
% results/main/rcaeval_re1/amber_zenodo_v2/summary_service.json
% results/analysis/final_ar_diagnostics/summary.json
% results/analysis/final_ar_diagnostics/case_diagnostics.csv
% notes/experiments/2026-08-28_final_counterfactual_ar_brushup.md
%
% E4：既存方法との比較（関連研究・問い1）
% notes/decisions/2026-08-29_metric_baseline_protocols.md
% notes/decisions/2026-08-29_baro_baseline_protocol.md
% notes/experiments/2026-08-31_run_aggregation_and_baseline_results.md
% results/baselines/rcaeval_re1/
% 各実験のsummary_service.jsonとcase JSONで確認する。
% BARO（正常時からのずれを得点にする比較手法）は開始時刻既知のRobustScorerのみ。
% ε-Diagnosis・RCD・CIRCA・RUNは比較手法名。本文導入時に考え方を説明する。
% 前回確認時にRUNは未完了。CIRCA/RUN等の変更点と入力条件を比較表で開示する。
% 実時間とケース時間合計、並列数・実行方式を混同しない。
%
% E5：実測値を使い続けるARの診断（提案の動機・問い2）
% notes/experiments/2026-08-22_rcaeval_re1_zenodo_v2_formal_ablation.md
% results/analysis/ar_rank_movement/summary.json
% results/analysis/ar_signal_attenuation/case_diagnostics.csv
% notes/ideas/2026-08-18_observed_lag_ar_signal_attenuation.md
% 旧Observed-lag AR/Raw × BF/GLRTの比較と、現在の構成を混同しない。
% 前処理や標準化等の条件差を監査し、一つの要素の効果と断定しない。
%
% E6：設定への感度・次数の選択（問い3）
% results/analysis/final_ar_sensitivity/summary.md
% notes/experiments/2026-08-29_extended_ar_order_sensitivity.md
% notes/decisions/2026-08-29_ar_order_selection.md
% results/sensitivity/rcaeval_re1/
% AR(7)は探索的候補で未昇格。draft2までに次数を確定する。
%
% E7：単位変換・正常時の得点（問い3・考察）
% results/analysis/ar_unit_invariance_ulp_guard/summary.json
% results/analysis/ar_final_sensitivity_pseudo_fault/unit_invariant_r0_98_p3/summary.json
% results/analysis/ar_residual_generalization_unit_invariant/summary.md
% 順位の一致と、得点・係数の数値的一致は別に確認する。
% 正常データの分割診断は、学習内分散の過小評価だけで問題を説明する説を支持しない。
%
% E8：採用しなかった案（考察・補足資料）
% notes/decisions/2026-08-27_bsrc_ar_research_stop.md
% notes/experiments/2026-08-25_bsrc_ar_adaptive_variance.md
% notes/experiments/2026-08-24_adaptive_direct_rollback_ablation.md
% BSRC-ARはAR生成過程の変化を直接比較する試行名で、主方式ではない。
% 層化75ケースの確認と正式375ケースの結果は分ける。執筆中に再開発しない。
